
\documentclass{article}
\usepackage{colortbl}
\usepackage{makecell}
\usepackage{multirow}
\usepackage{supertabular}

\begin{document}

\newcounter{utterance}

\centering \large Interaction Transcript for game `function\_detective', experiment `two\_numbers\_passive\_examples', episode 289 with Qwen3.6{-}35B{-}A3B{-}FP8{-}without{-}reasoning.
\vspace{24pt}

{ \footnotesize  \setcounter{utterance}{1}
\setlength{\tabcolsep}{0pt}
\begin{supertabular}{c@{$\;$}|p{.15\linewidth}@{}p{.15\linewidth}p{.15\linewidth}p{.15\linewidth}p{.15\linewidth}p{.15\linewidth}}
    \# & \multicolumn{2}{c}{Player} && \multicolumn{2}{c}{Game Master} \\
    \hline

    \theutterance \stepcounter{utterance}  
    & & & \multicolumn{4}{p{0.6\linewidth}}{
        \cellcolor[rgb]{0.9,0.9,0.9}{
            \makecell[{{p{\linewidth}}}]{
                \texttt{\tiny{[P1$\langle$GM]}}
                \texttt{I have a hidden Python function for you to figure out.} \\
\\ 
\texttt{Your job is to infer what it does and then submit a correct Python implementation.} \\
\\ 
\texttt{Function signature:} \\
\texttt{(a: int, b: int) {-}> int} \\
\\ 
\texttt{You have 15 turns total.} \\
\\ 
\texttt{I control which valid input{-}output examples you see. You can request them one at a time.} \\
\\ 
\texttt{You must make your best SOLVE attempt no later than the second{-}to{-}last turn.} \\
\texttt{Do not wait until the very end.} \\
\\ 
\texttt{You can use these actions:} \\
\texttt{{-} JSON format described at the end of the prompt} \\
\texttt{{-} SOLVE: ```python ... ```} \\
\\ 
\texttt{Do not add import statements in your guess.} \\
\texttt{Anything you might need is already available, so just write the function body in hypothesis and return the answer.} \\
\\ 
\texttt{Preloaded examples:} \\
\texttt{1. (({-}99, 10), 723)} \\
\\ 
\texttt{Scoring:} \\
\texttt{{-} correctness is required for any non{-}zero score} \\
\texttt{{-} in interactive modes, solving with fewer information requests is better} \\
\\ 
\texttt{Input formatting reminders:} \\
\texttt{{-} strings must be quoted} \\
\texttt{{-} booleans must be True/False} \\
\texttt{{-} use commas between multiple values} \\
\texttt{{-} do not write any extra action tags in the same message} \\
\\ 
\texttt{Format example:} \\
\texttt{You write:} \\
\texttt{\{'current\_mode': <probing/confirming/don't know; choose whichever best describes your current state in the task>, 'input\_rationale': <brief description of why you chose this hypothesis>, 'hypothesis': <either the exact string "unknown" if you have no guess yet OR your current best guess as runnable Python matching the given signature, with a one{-}line docstring stating what it does. Never write placeholder code, if you cannot commit to a guess, use "unknown". Match the given signature.} \\
\texttt{e.g.} \\
\texttt{def f(x):} \\
\texttt{	"""Returns the absolute value of x."""} \\
\texttt{	return x if x >= 0 else {-}x} \\
\\ 
\texttt{Use "unknown" if you have no guess.>, 'confidence\_score': <score between 1{-}100. 100 means highest confidence in your hypothesis>\}} \\
\\ 
\texttt{Output only one valid json per turn.} \\
            }
        }
    }
     \\ \\

    \theutterance \stepcounter{utterance}  
    & \multicolumn{4}{p{0.6\linewidth}}{
        \cellcolor[rgb]{0.561,0.737,0.561}{
            \makecell[{{p{\linewidth}}}]{
                \texttt{\tiny{[P1$\rangle$GM]}}
                \texttt{\{'current\_mode': 'probing', 'input\_rationale': 'I have one example: f({-}99, 10) = 723. I need more examples to understand the relationship between a, b, and the result. I will request another example with different values to gather more data.', 'hypothesis': 'unknown', 'confidence\_score': 10\}} \\
            }
        }
    }
    & & \\ \\

    \theutterance \stepcounter{utterance}  
    & & & \multicolumn{4}{p{0.6\linewidth}}{
        \cellcolor[rgb]{0.9,0.9,0.9}{
            \makecell[{{p{\linewidth}}}]{
                \texttt{\tiny{[P1$\langle$GM]}}
                \texttt{Game Violation: ParseError: Response is not valid JSON.} \\
\texttt{Output ONLY:} \\
\texttt{One JSON with confidence\_score, current\_mode, input\_rationale, and hypothesis keys ...} \\
\texttt{OR} \\
\texttt{SOLVE: ```python ... ```} \\
            }
        }
    }
     \\ \\

    \theutterance \stepcounter{utterance}  
    & \multicolumn{4}{p{0.6\linewidth}}{
        \cellcolor[rgb]{0.561,0.737,0.561}{
            \makecell[{{p{\linewidth}}}]{
                \texttt{\tiny{[P1$\rangle$GM]}}
                \texttt{\{'current\_mode': 'probing', 'input\_rationale': 'I have one example: f({-}99, 10) = 723. I need more examples to understand the relationship between a, b, and the result. I will request another example with different values to gather more data.', 'hypothesis': 'unknown', 'confidence\_score': 10\}} \\
            }
        }
    }
    & & \\ \\

    \theutterance \stepcounter{utterance}  
    & & & \multicolumn{4}{p{0.6\linewidth}}{
        \cellcolor[rgb]{0.9,0.9,0.9}{
            \makecell[{{p{\linewidth}}}]{
                \texttt{\tiny{[P1$\langle$GM]}}
                \texttt{Game Violation: ParseError: Response is not valid JSON.} \\
\texttt{Output ONLY:} \\
\texttt{One JSON with confidence\_score, current\_mode, input\_rationale, and hypothesis keys ...} \\
\texttt{OR} \\
\texttt{SOLVE: ```python ... ```} \\
            }
        }
    }
     \\ \\

    \theutterance \stepcounter{utterance}  
    & \multicolumn{4}{p{0.6\linewidth}}{
        \cellcolor[rgb]{0.561,0.737,0.561}{
            \makecell[{{p{\linewidth}}}]{
                \texttt{\tiny{[P1$\rangle$GM]}}
                \texttt{\{"current\_mode": "probing", "input\_rationale": "I have one example: f({-}99, 10) = 723. I need more examples to understand the relationship between a, b, and the result. I will request another example with different values to gather more data.", "hypothesis": "unknown", "confidence\_score": 10\}} \\
            }
        }
    }
    & & \\ \\

    \theutterance \stepcounter{utterance}  
    & & & \multicolumn{4}{p{0.6\linewidth}}{
        \cellcolor[rgb]{0.9,0.9,0.9}{
            \makecell[{{p{\linewidth}}}]{
                \texttt{\tiny{[P1$\langle$GM]}}
                \texttt{OUTPUT: (({-}1, {-}99), {-}399)} \\
            }
        }
    }
     \\ \\

    \theutterance \stepcounter{utterance}  
    & \multicolumn{4}{p{0.6\linewidth}}{
        \cellcolor[rgb]{0.561,0.737,0.561}{
            \makecell[{{p{\linewidth}}}]{
                \texttt{\tiny{[P1$\rangle$GM]}}
                \texttt{\{"current\_mode": "probing", "input\_rationale": "I have two examples: f({-}99, 10) = 723 and f({-}1, {-}99) = {-}399. Let's analyze the relationship. \textbackslash{}n\textbackslash{}nExample 1: a={-}99, b=10, result=723.\textbackslash{}nExample 2: a={-}1, b={-}99, result={-}399.\textbackslash{}n\textbackslash{}nLet's try simple linear combinations: c1*a + c2*b = result.\textbackslash{}n1) {-}99*c1 + 10*c2 = 723\textbackslash{}n2) {-}1*c1 {-} 99*c2 = {-}399 => c1 = 399 {-} 99*c2\textbackslash{}n\textbackslash{}nSubstitute c1 into 1):\textbackslash{}n{-}99*(399 {-} 99*c2) + 10*c2 = 723\textbackslash{}n{-}39501 + 9801*c2 + 10*c2 = 723\textbackslash{}n9811*c2 = 40224\textbackslash{}nc2 = 40224 / 9811 ≈ 4.1\textbackslash{}nThis doesn't look like integer coefficients. Let's try other operations.\textbackslash{}n\textbackslash{}nMaybe it involves multiplication? \textbackslash{}n{-}99 * 10 = {-}990} \\
            }
        }
    }
    & & \\ \\

    \theutterance \stepcounter{utterance}  
    & & & \multicolumn{4}{p{0.6\linewidth}}{
        \cellcolor[rgb]{0.9,0.9,0.9}{
            \makecell[{{p{\linewidth}}}]{
                \texttt{\tiny{[P1$\langle$GM]}}
                \texttt{Game Violation: ParseError: Response is not valid JSON.} \\
\texttt{Output ONLY:} \\
\texttt{One JSON with confidence\_score, current\_mode, input\_rationale, and hypothesis keys ...} \\
\texttt{OR} \\
\texttt{SOLVE: ```python ... ```} \\
            }
        }
    }
     \\ \\

    \theutterance \stepcounter{utterance}  
    & \multicolumn{4}{p{0.6\linewidth}}{
        \cellcolor[rgb]{0.561,0.737,0.561}{
            \makecell[{{p{\linewidth}}}]{
                \texttt{\tiny{[P1$\rangle$GM]}}
                \texttt{\{"current\_mode": "probing", "input\_rationale": "I have two examples: f({-}99, 10) = 723 and f({-}1, {-}99) = {-}399. I need a third example to better constrain the possible functions, especially to distinguish between linear combinations, multiplicative terms, or other operations.", "hypothesis": "unknown", "confidence\_score": 10\}} \\
            }
        }
    }
    & & \\ \\

    \theutterance \stepcounter{utterance}  
    & & & \multicolumn{4}{p{0.6\linewidth}}{
        \cellcolor[rgb]{0.9,0.9,0.9}{
            \makecell[{{p{\linewidth}}}]{
                \texttt{\tiny{[P1$\langle$GM]}}
                \texttt{OUTPUT: (({-}99, {-}2), {-}501)} \\
            }
        }
    }
     \\ \\

    \theutterance \stepcounter{utterance}  
    & \multicolumn{4}{p{0.6\linewidth}}{
        \cellcolor[rgb]{0.561,0.737,0.561}{
            \makecell[{{p{\linewidth}}}]{
                \texttt{\tiny{[P1$\rangle$GM]}}
                \texttt{\{"current\_mode": "confirming", "input\_rationale": "I have three examples: f({-}99, 10) = 723, f({-}1, {-}99) = {-}399, f({-}99, {-}2) = {-}501. Let's test a linear hypothesis: f(a, b) = c1*a + c2*b. From example 1: {-}99*c1 + 10*c2 = 723. From example 3: {-}99*c1 {-} 2*c2 = {-}501. Subtracting the two equations: 12*c2 = 1224 => c2 = 102. Then {-}99*c1 + 1020 = 723 => {-}99*c1 = {-}297 => c1 = 3. So the hypothesis is f(a, b) = 3*a + 102*b. Let's check with example 2: 3*({-}1) + 102*({-}99) = {-}3 {-} 10098 = {-}10101. This does not match {-}399. So it's not a simple linear combination. Let's re{-}examine. Maybe it's f(a, b) = a * b + something? {-}99*10 = {-}990. 7} \\
            }
        }
    }
    & & \\ \\

    \theutterance \stepcounter{utterance}  
    & & & \multicolumn{4}{p{0.6\linewidth}}{
        \cellcolor[rgb]{0.9,0.9,0.9}{
            \makecell[{{p{\linewidth}}}]{
                \texttt{\tiny{[P1$\langle$GM]}}
                \texttt{Game Violation: ParseError: Response is not valid JSON.} \\
\texttt{Output ONLY:} \\
\texttt{One JSON with confidence\_score, current\_mode, input\_rationale, and hypothesis keys ...} \\
\texttt{OR} \\
\texttt{SOLVE: ```python ... ```} \\
            }
        }
    }
     \\ \\

    \theutterance \stepcounter{utterance}  
    & \multicolumn{4}{p{0.6\linewidth}}{
        \cellcolor[rgb]{0.561,0.737,0.561}{
            \makecell[{{p{\linewidth}}}]{
                \texttt{\tiny{[P1$\rangle$GM]}}
                \texttt{\{"current\_mode": "confirming", "input\_rationale": "I have three examples: f({-}99, 10) = 723, f({-}1, {-}99) = {-}399, f({-}99, {-}2) = {-}501. Let's test a linear hypothesis: f(a, b) = c1*a + c2*b. From example 1: {-}99*c1 + 10*c2 = 723. From example 3: {-}99*c1 {-} 2*c2 = {-}501. Subtracting the two equations: 12*c2 = 1224 => c2 = 102. Then {-}99*c1 + 1020 = 723 => {-}99*c1 = {-}297 => c1 = 3. So the hypothesis is f(a, b) = 3*a + 102*b. Let's check with example 2: 3*({-}1) + 102*({-}99) = {-}3 {-} 10098 = {-}10101. This does not match {-}399. So it's not a simple linear combination. Let's re{-}examine. Maybe it's f(a, b) = a * b + something? {-}99*10 = {-}990. 7} \\
            }
        }
    }
    & & \\ \\

    \theutterance \stepcounter{utterance}  
    & & & \multicolumn{4}{p{0.6\linewidth}}{
        \cellcolor[rgb]{0.9,0.9,0.9}{
            \makecell[{{p{\linewidth}}}]{
                \texttt{\tiny{[P1$\langle$GM]}}
                \texttt{Game Violation: ParseError: Response is not valid JSON.} \\
\texttt{Output ONLY:} \\
\texttt{One JSON with confidence\_score, current\_mode, input\_rationale, and hypothesis keys ...} \\
\texttt{OR} \\
\texttt{SOLVE: ```python ... ```} \\
            }
        }
    }
     \\ \\

    \theutterance \stepcounter{utterance}  
    & \multicolumn{4}{p{0.6\linewidth}}{
        \cellcolor[rgb]{0.561,0.737,0.561}{
            \makecell[{{p{\linewidth}}}]{
                \texttt{\tiny{[P1$\rangle$GM]}}
                \texttt{\{"current\_mode": "confirming", "input\_rationale": "I have three examples: f({-}99, 10) = 723, f({-}1, {-}99) = {-}399, f({-}99, {-}2) = {-}501. Let's test a linear hypothesis: f(a, b) = c1*a + c2*b. From example 1: {-}99*c1 + 10*c2 = 723. From example 3: {-}99*c1 {-} 2*c2 = {-}501. Subtracting the two equations: 12*c2 = 1224 => c2 = 102. Then {-}99*c1 + 1020 = 723 => {-}99*c1 = {-}297 => c1 = 3. So the hypothesis is f(a, b) = 3*a + 102*b. Let's check with example 2: 3*({-}1) + 102*({-}99) = {-}3 {-} 10098 = {-}10101. This does not match {-}399. So it's not a simple linear combination. Let's re{-}examine. Maybe it's f(a, b) = a * b + something? {-}99*10 = {-}990. 7} \\
            }
        }
    }
    & & \\ \\

    \theutterance \stepcounter{utterance}  
    & & & \multicolumn{4}{p{0.6\linewidth}}{
        \cellcolor[rgb]{0.9,0.9,0.9}{
            \makecell[{{p{\linewidth}}}]{
                \texttt{\tiny{[P1$\langle$GM]}}
                \texttt{Game Violation: ParseError: Response is not valid JSON.} \\
\texttt{Output ONLY:} \\
\texttt{One JSON with confidence\_score, current\_mode, input\_rationale, and hypothesis keys ...} \\
\texttt{OR} \\
\texttt{SOLVE: ```python ... ```} \\
            }
        }
    }
     \\ \\

    \theutterance \stepcounter{utterance}  
    & \multicolumn{4}{p{0.6\linewidth}}{
        \cellcolor[rgb]{0.561,0.737,0.561}{
            \makecell[{{p{\linewidth}}}]{
                \texttt{\tiny{[P1$\rangle$GM]}}
                \texttt{\{"current\_mode": "confirming", "input\_rationale": "I have three examples: f({-}99, 10) = 723, f({-}1, {-}99) = {-}399, f({-}99, {-}2) = {-}501. Let's test a linear hypothesis: f(a, b) = c1*a + c2*b. From example 1: {-}99*c1 + 10*c2 = 723. From example 3: {-}99*c1 {-} 2*c2 = {-}501. Subtracting the two equations: 12*c2 = 1224 => c2 = 102. Then {-}99*c1 + 1020 = 723 => {-}99*c1 = {-}297 => c1 = 3. So the hypothesis is f(a, b) = 3*a + 102*b. Let's check with example 2: 3*({-}1) + 102*({-}99) = {-}3 {-} 10098 = {-}10101. This does not match {-}399. So it's not a simple linear combination. Let's re{-}examine. Maybe it's f(a, b) = a * b + something? {-}99*10 = {-}990. 7} \\
            }
        }
    }
    & & \\ \\

    \theutterance \stepcounter{utterance}  
    & & & \multicolumn{4}{p{0.6\linewidth}}{
        \cellcolor[rgb]{0.9,0.9,0.9}{
            \makecell[{{p{\linewidth}}}]{
                \texttt{\tiny{[P1$\langle$GM]}}
                \texttt{Game Violation: ParseError: Response is not valid JSON.} \\
\texttt{Output ONLY:} \\
\texttt{One JSON with confidence\_score, current\_mode, input\_rationale, and hypothesis keys ...} \\
\texttt{OR} \\
\texttt{SOLVE: ```python ... ```} \\
            }
        }
    }
     \\ \\

    \theutterance \stepcounter{utterance}  
    & \multicolumn{4}{p{0.6\linewidth}}{
        \cellcolor[rgb]{0.561,0.737,0.561}{
            \makecell[{{p{\linewidth}}}]{
                \texttt{\tiny{[P1$\rangle$GM]}}
                \texttt{\{"current\_mode": "confirming", "input\_rationale": "I have three examples: f({-}99, 10) = 723, f({-}1, {-}99) = {-}399, f({-}99, {-}2) = {-}501. Let's test a linear hypothesis: f(a, b) = c1*a + c2*b. From example 1: {-}99*c1 + 10*c2 = 723. From example 3: {-}99*c1 {-} 2*c2 = {-}501. Subtracting the two equations: 12*c2 = 1224 => c2 = 102. Then {-}99*c1 + 1020 = 723 => {-}99*c1 = {-}297 => c1 = 3. So the hypothesis is f(a, b) = 3*a + 102*b. Let's check with example 2: 3*({-}1) + 102*({-}99) = {-}3 {-} 10098 = {-}10101. This does not match {-}399. So it's not a simple linear combination. Let's re{-}examine. Maybe it's f(a, b) = a * b + something? {-}99*10 = {-}990. 7} \\
            }
        }
    }
    & & \\ \\

    \theutterance \stepcounter{utterance}  
    & & & \multicolumn{4}{p{0.6\linewidth}}{
        \cellcolor[rgb]{0.9,0.9,0.9}{
            \makecell[{{p{\linewidth}}}]{
                \texttt{\tiny{[P1$\langle$GM]}}
                \texttt{Game Violation: ParseError: Response is not valid JSON.} \\
\texttt{Output ONLY:} \\
\texttt{One JSON with confidence\_score, current\_mode, input\_rationale, and hypothesis keys ...} \\
\texttt{OR} \\
\texttt{SOLVE: ```python ... ```} \\
            }
        }
    }
     \\ \\

    \theutterance \stepcounter{utterance}  
    & \multicolumn{4}{p{0.6\linewidth}}{
        \cellcolor[rgb]{0.561,0.737,0.561}{
            \makecell[{{p{\linewidth}}}]{
                \texttt{\tiny{[P1$\rangle$GM]}}
                \texttt{\{"current\_mode": "confirming", "input\_rationale": "I have three examples: f({-}99, 10) = 723, f({-}1, {-}99) = {-}399, f({-}99, {-}2) = {-}501. Let's test a linear hypothesis: f(a, b) = c1*a + c2*b. From example 1: {-}99*c1 + 10*c2 = 723. From example 3: {-}99*c1 {-} 2*c2 = {-}501. Subtracting the two equations: 12*c2 = 1224 => c2 = 102. Then {-}99*c1 + 1020 = 723 => {-}99*c1 = {-}297 => c1 = 3. So the hypothesis is f(a, b) = 3*a + 102*b. Let's check with example 2: 3*({-}1) + 102*({-}99) = {-}3 {-} 10098 = {-}10101. This does not match {-}399. So it's not a simple linear combination. Let's re{-}examine. Maybe it's f(a, b) = a * b + something? {-}99*10 = {-}990. 7} \\
            }
        }
    }
    & & \\ \\

    \theutterance \stepcounter{utterance}  
    & & & \multicolumn{4}{p{0.6\linewidth}}{
        \cellcolor[rgb]{0.9,0.9,0.9}{
            \makecell[{{p{\linewidth}}}]{
                \texttt{\tiny{[P1$\langle$GM]}}
                \texttt{Game Violation: ParseError: Response is not valid JSON.} \\
\texttt{Output ONLY:} \\
\texttt{One JSON with confidence\_score, current\_mode, input\_rationale, and hypothesis keys ...} \\
\texttt{OR} \\
\texttt{SOLVE: ```python ... ```} \\
            }
        }
    }
     \\ \\

    \theutterance \stepcounter{utterance}  
    & \multicolumn{4}{p{0.6\linewidth}}{
        \cellcolor[rgb]{0.561,0.737,0.561}{
            \makecell[{{p{\linewidth}}}]{
                \texttt{\tiny{[P1$\rangle$GM]}}
                \texttt{\{"current\_mode": "confirming", "input\_rationale": "I have three examples: f({-}99, 10) = 723, f({-}1, {-}99) = {-}399, f({-}99, {-}2) = {-}501. Let's test a linear hypothesis: f(a, b) = c1*a + c2*b. From example 1: {-}99*c1 + 10*c2 = 723. From example 3: {-}99*c1 {-} 2*c2 = {-}501. Subtracting the two equations: 12*c2 = 1224 => c2 = 102. Then {-}99*c1 + 1020 = 723 => {-}99*c1 = {-}297 => c1 = 3. So the hypothesis is f(a, b) = 3*a + 102*b. Let's check with example 2: 3*({-}1) + 102*({-}99) = {-}3 {-} 10098 = {-}10101. This does not match {-}399. So it's not a simple linear combination. Let's re{-}examine. Maybe it's f(a, b) = a * b + something? {-}99*10 = {-}990. 7} \\
            }
        }
    }
    & & \\ \\

    \theutterance \stepcounter{utterance}  
    & & & \multicolumn{4}{p{0.6\linewidth}}{
        \cellcolor[rgb]{0.9,0.9,0.9}{
            \makecell[{{p{\linewidth}}}]{
                \texttt{\tiny{[P1$\langle$GM]}}
                \texttt{Game Violation: ParseError: Response is not valid JSON.} \\
\texttt{Output ONLY:} \\
\texttt{One JSON with confidence\_score, current\_mode, input\_rationale, and hypothesis keys ...} \\
\texttt{OR} \\
\texttt{SOLVE: ```python ... ```} \\
            }
        }
    }
     \\ \\

    \theutterance \stepcounter{utterance}  
    & \multicolumn{4}{p{0.6\linewidth}}{
        \cellcolor[rgb]{0.561,0.737,0.561}{
            \makecell[{{p{\linewidth}}}]{
                \texttt{\tiny{[P1$\rangle$GM]}}
                \texttt{\{"current\_mode": "confirming", "input\_rationale": "I have three examples: f({-}99, 10) = 723, f({-}1, {-}99) = {-}399, f({-}99, {-}2) = {-}501. Let's test a linear hypothesis: f(a, b) = c1*a + c2*b. From example 1: {-}99*c1 + 10*c2 = 723. From example 3: {-}99*c1 {-} 2*c2 = {-}501. Subtracting the two equations: 12*c2 = 1224 => c2 = 102. Then {-}99*c1 + 1020 = 723 => {-}99*c1 = {-}297 => c1 = 3. So the hypothesis is f(a, b) = 3*a + 102*b. Let's check with example 2: 3*({-}1) + 102*({-}99) = {-}3 {-} 10098 = {-}10101. This does not match {-}399. So it's not a simple linear combination. Let's re{-}examine. Maybe it's f(a, b) = a * b + something? {-}99*10 = {-}990. 7} \\
            }
        }
    }
    & & \\ \\

    \theutterance \stepcounter{utterance}  
    & & & \multicolumn{4}{p{0.6\linewidth}}{
        \cellcolor[rgb]{0.9,0.9,0.9}{
            \makecell[{{p{\linewidth}}}]{
                \texttt{\tiny{[P1$\langle$GM]}}
                \texttt{Game Violation: ParseError: Response is not valid JSON.} \\
\texttt{Output ONLY:} \\
\texttt{One JSON with confidence\_score, current\_mode, input\_rationale, and hypothesis keys ...} \\
\texttt{OR} \\
\texttt{SOLVE: ```python ... ```} \\
            }
        }
    }
     \\ \\

    \theutterance \stepcounter{utterance}  
    & \multicolumn{4}{p{0.6\linewidth}}{
        \cellcolor[rgb]{0.561,0.737,0.561}{
            \makecell[{{p{\linewidth}}}]{
                \texttt{\tiny{[P1$\rangle$GM]}}
                \texttt{\{"current\_mode": "confirming", "input\_rationale": "I have three examples: f({-}99, 10) = 723, f({-}1, {-}99) = {-}399, f({-}99, {-}2) = {-}501. Let's test a linear hypothesis: f(a, b) = c1*a + c2*b. From example 1: {-}99*c1 + 10*c2 = 723. From example 3: {-}99*c1 {-} 2*c2 = {-}501. Subtracting the two equations: 12*c2 = 1224 => c2 = 102. Then {-}99*c1 + 1020 = 723 => {-}99*c1 = {-}297 => c1 = 3. So the hypothesis is f(a, b) = 3*a + 102*b. Let's check with example 2: 3*({-}1) + 102*({-}99) = {-}3 {-} 10098 = {-}10101. This does not match {-}399. So it's not a simple linear combination. Let's re{-}examine. Maybe it's f(a, b) = a * b + something? {-}99*10 = {-}990. 7} \\
            }
        }
    }
    & & \\ \\

    \theutterance \stepcounter{utterance}  
    & & & \multicolumn{4}{p{0.6\linewidth}}{
        \cellcolor[rgb]{0.9,0.9,0.9}{
            \makecell[{{p{\linewidth}}}]{
                \texttt{\tiny{[P1$\langle$GM]}}
                \texttt{Game Violation: ParseError: Response is not valid JSON.} \\
\texttt{Output ONLY:} \\
\texttt{One JSON with confidence\_score, current\_mode, input\_rationale, and hypothesis keys ...} \\
\texttt{OR} \\
\texttt{SOLVE: ```python ... ```} \\
            }
        }
    }
     \\ \\

    \theutterance \stepcounter{utterance}  
    & \multicolumn{4}{p{0.6\linewidth}}{
        \cellcolor[rgb]{0.561,0.737,0.561}{
            \makecell[{{p{\linewidth}}}]{
                \texttt{\tiny{[P1$\rangle$GM]}}
                \texttt{\{"current\_mode": "confirming", "input\_rationale": "I have three examples: f({-}99, 10) = 723, f({-}1, {-}99) = {-}399, f({-}99, {-}2) = {-}501. Let's test a linear hypothesis: f(a, b) = c1*a + c2*b. From example 1: {-}99*c1 + 10*c2 = 723. From example 3: {-}99*c1 {-} 2*c2 = {-}501. Subtracting the two equations: 12*c2 = 1224 => c2 = 102. Then {-}99*c1 + 1020 = 723 => {-}99*c1 = {-}297 => c1 = 3. So the hypothesis is f(a, b) = 3*a + 102*b. Let's check with example 2: 3*({-}1) + 102*({-}99) = {-}3 {-} 10098 = {-}10101. This does not match {-}399. So it's not a simple linear combination. Let's re{-}examine. Maybe it's f(a, b) = a * b + something? {-}99*10 = {-}990. 7} \\
            }
        }
    }
    & & \\ \\

    \theutterance \stepcounter{utterance}  
    & & & \multicolumn{2}{p{0.3\linewidth}}{
        \cellcolor[rgb]{0.9,0.9,0.9}{
            \makecell[{{p{\linewidth}}}]{
                \texttt{\tiny{[GM$|$GM]}}
                \texttt{Game Verdict: LOSS (turn limit reached)} \\
            }
        }
    }
    & & \\ \\

\end{supertabular}
}

\end{document}
